\section{Introduction}
A signal-reconnaissance system combines reports of detected activity from several RF receivers. A recognizer may also supply probabilities over signal families \cite{syed2023,zhang2022amc,tian2026}. Differences in units, class orders, and measurement errors create multi-source heterogeneity. All sources in this study are RF receivers.

Here, \emph{semantic evidence} is probability compatibility after class columns have been mapped to common meanings; it does not refer to natural-language interpretation or decoded message content. For example, column 1 can denote different families at two receivers, so comparing those columns directly can produce false agreement. Correct mapping makes the probabilities comparable, but separate transmissions can still share a family. Event identity therefore requires a further inference step.

We distinguish three alignment levels. Representation alignment converts values and uncertainties to common units. Class-semantic alignment registers native probability columns to a shared signal-family vocabulary. Instance alignment groups reports of one bounded transmission. The title's \emph{consistency} refers to the resulting groups: membership is transitive, each receiver appears at most once per group, and every pair passes the physical gate. These requirements enforce the input contract without certifying event identity.

\widefig{fig01_task}{The alignment task. Native units and class orders differ across RF receivers. Hidden source positions illustrate the generating scene. They are not estimated or supplied to inference. A receiver-exclusive partition is a task requirement. It is not an observed accuracy result. A valid group also requires every pair to pass physical admission. The panels are conceptual examples.}{fig:motivation}

The source registry states field meanings, units, and uncertainty declarations \cite{ssn,sosa,emeterio2026}. The baseline also supplies class maps, while the extension estimates selected missing permutations within a known four-class vocabulary. One reference receiver retains its known map to fix the otherwise arbitrary common class naming. The method does not learn field definitions or new classes.

Unit conversion, calibrated scoring, assignment, and signed clustering are established components \cite{guo,kuhn,bansal,baumann2025}. Permutation synchronization also estimates several maps jointly \cite{pachauri2013,chen2014}. Here, class correspondences must be recovered from uncertain RF reports whose event groups are initially unknown. We address this report-level alignment problem and test its limits:
\begin{enumerate}[leftmargin=1.5em,itemsep=2pt]
\item We specify an auditable RF observation contract. It separates representation, class-semantic, and instance alignment. Matched unit and vocabulary re-encodings leave canonical evidence and event partitions unchanged.
\item We formulate missing class maps and event groups in one score model. Receiver-pair tables retain every candidate relative map. Exact conditional search provides a small-problem reference. We compare it with coordinate updates, reference-only fitting, and an adapted spectral synchronization method. A conditional identification argument states the role of a known reference and informative pair scores.
\item We test the alignment across frozen scorers and independent descriptor scenes. Same-class collisions and classifier-error interventions expose the limits of semantic evidence. Integer programs measure the cost of approximate event grouping. Public records provide a separate check of metadata correspondence and clock preprocessing.
\end{enumerate}

The inference target is a bounded transmission. Receiver locations are known, while source locations and event count are hidden. The method starts after detection and does not evaluate a waveform front end \cite{hannotier2024,kropfreiter2024,yan2024}.

Controlled descriptors isolate class-map and association effects. We retain the initial 800-scene benchmark for continuity. Separate adaptation and test scenes support the mapping comparisons and semantic stress tests. Public LoED records provide an auxiliary metadata check \cite{loed2020,loeddataset}. We use them to test reference-packet correspondence and preprocessing sensitivity. The controlled benchmark supplies the evidence for the class-probability mechanism.
\FloatBarrier
